\chapter{Radioactivity and Nuclear Reactions}
\label{chap:exo_ch03}

\begin{quicksummary}{Decay Laws, Kinetics, Activity Units \& Radiometric Dating}
\begin{enumerate}[leftmargin=1.4em,itemsep=3pt]
  \item \textbf{Soddy-Fajans Conservation Laws:}
  In any nuclear transformation:
  \[
  \sum Z_{\text{reactants}} = \sum Z_{\text{products}}, \qquad 
  \sum A_{\text{reactants}} = \sum A_{\text{products}}.
  \]
  \item \textbf{Fundamental Decay Modes:}
  \begin{itemize}[leftmargin=1.2em,itemsep=1pt]
    \item \textit{Alpha ($\alpha$):} $\ce{^{A}_{Z}X -> ^{A-4}_{Z-2}Y + ^{4}_{2}He}$.
    \item \textit{Beta Minus ($\beta^-$):} $\ce{^{A}_{Z}X -> ^{A}_{Z+1}Y + ^{0}_{-1}e + \bar{\nu}_e}$ \quad ($n \to p + e^- + \bar{\nu}_e$).
    \item \textit{Beta Plus ($\beta^+$):} $\ce{^{A}_{Z}X -> ^{A}_{Z-1}Y + ^{0}_{+1}e + \nu_e}$ \quad ($p \to n + e^+ + \nu_e$).
    \item \textit{Electron Capture (EC):} $\ce{^{A}_{Z}X + ^{0}_{-1}e -> ^{A}_{Z-1}Y + \nu_e}$.
    \item \textit{Gamma De-excitation ($\gamma$):} $\ce{^{A}_{Z}X^* -> ^{A}_{Z}X + \gamma}$.
  \end{itemize}
  \item \textbf{Decay Kinetics \& Exponential Law:}
  \[
  \frac{\dd N}{\dd t} = -\lambda N(t) \implies N(t) = N_0 e^{-\lambda t}, \qquad m(t) = m_0 e^{-\lambda t}.
  \]
  \item \textbf{Half-Life ($t_{1/2}$) \& Mean Lifetime ($\tau$):}
  \[
  t_{1/2} = \frac{\ln 2}{\lambda} = \frac{0.693147}{\lambda}, \qquad \tau = \frac{1}{\lambda} = \frac{t_{1/2}}{\ln 2}.
  \]
  After $k$ half-lives ($t = k \cdot t_{1/2}$): $N(t) = \dfrac{N_0}{2^k}$.
  \item \textbf{Activity ($A$) \& Units:}
  \[
  A(t) = -\frac{\dd N}{\dd t} = \lambda N(t) = A_0 e^{-\lambda t}.
  \]
  \[
  1\ \mathrm{Becquerel\ (Bq)} = 1\ \mathrm{decay\cdot s^{-1}}, \qquad 1\ \mathrm{Curie\ (Ci)} = 3.7 \times 10^{10}\ \mathrm{Bq}.
  \]
  \item \textbf{Radiometric Dating Formulas:}
  \begin{itemize}[leftmargin=1.2em,itemsep=1pt]
    \item \textit{Radiocarbon ($^{14}\ce{C}$):} $t = \dfrac{t_{1/2}}{\ln 2} \ln \left( \dfrac{A_0}{A(t)} \right) = \dfrac{5730\ \mathrm{yr}}{\ln 2} \ln \left( \dfrac{A_0}{A(t)} \right)$.
    \item \textit{Uranium-Lead ($^{238}\ce{U} \to ^{206}\ce{Pb}$):} $t = \dfrac{1}{\lambda} \ln \left( 1 + \dfrac{N(\ce{^{206}Pb})}{N(\ce{^{238}U})} \right)$.
  \end{itemize}
\end{enumerate}
\end{quicksummary}

\vspace{10pt}

% ===========================================================================
% EXERCISE 3.1
% ===========================================================================
\begin{exercise}[$\bigstar$ Foundational]{Physical Properties and Shielding of Nuclear Radiations}
Summarize the fundamental physical properties of the four primary ionizing radiations:
\begin{enumerate}[leftmargin=1.5em,itemsep=2pt]
  \item Complete a comparative table specifying for $\alpha$, $\beta^-$, $\beta^+$, and $\gamma$:
  \begin{itemize}[itemsep=1pt]
    \item Particle nature and composition;
    \item Electric charge (in units of elementary charge $e$);
    \item Rest mass (in atomic mass units $\mathrm{u}$);
    \item Penetrating power and typical shielding material required to stop it.
  \end{itemize}
  \item Explain why alpha radiation is benign outside the human body but exceptionally lethal when inhaled or ingested.
\end{enumerate}
\end{exercise}

\begin{tabletworkspace}[5.0cm]
\end{tabletworkspace}

\begin{solution}
\textbf{1. Comparative Table of Radiations:}
\begin{center}
\renewcommand{\arraystretch}{1.3}
\begin{tabular}{@{}p{2.3cm}p{3.1cm}p{1.8cm}p{2.3cm}p{5.5cm}@{}}
\toprule
\textbf{Radiation} & \textbf{Nature} & \textbf{Charge} & \textbf{Rest Mass} & \textbf{Shielding Required} \\
\midrule
Alpha ($\alpha$) & $\ce{^{4}_{2}He}^{2+}$ nucleus & $+2\,e$ & $4.0015\ \mathrm{u}$ & Sheet of paper, dead skin layer ($40\ \mu\mathrm{m}$) \\
Beta Minus ($\beta^-$) & Fast electron ($e^-$) & $-1\,e$ & $0.000549\ \mathrm{u}$ & Few mm of aluminum or Plexiglas \\
Beta Plus ($\beta^+$) & Positron ($e^+$) & $+1\,e$ & $0.000549\ \mathrm{u}$ & Annihilates with $e^-$ into two $511\ \mathrm{keV}\ \gamma$ photons \\
Gamma ($\gamma$) & High-energy photon & $0$ & $0$ & Several cm of lead ($\ce{Pb}$) or dense concrete \\
\bottomrule
\end{tabular}
\end{center}

\textbf{2. Internal vs. External Biological Hazard of $\alpha$-Radiation:}
\begin{itemize}[leftmargin=1.5em,itemsep=2pt]
  \item \textbf{Externally:} Alpha particles have very low penetration ($< 0.05\ \mathrm{mm}$ in tissue). They are completely stopped by the dead stratum corneum layer of human skin and cannot reach living cells.
  \item \textbf{Internally (Inhaled/Ingested):} Alpha particles possess high mass and $+2e$ charge, resulting in an extremely high \textbf{linear energy transfer (LET)}. They dump their full kinetic energy ($\approx 4\text{--}8\ \mathrm{MeV}$) within a tiny cellular volume ($\approx 40\ \mu\mathrm{m}$, the diameter of 2--3 cells), causing severe double-strand DNA breaks, chromosomal aberrations, and cellular death.
\end{itemize}
\end{solution}

% ===========================================================================
% EXERCISE 3.2
% ===========================================================================
\begin{exercise}[$\bigstar$ Foundational]{Soddy-Fajans Conservation Laws in Alpha Decay}
Applying the Soddy-Fajans conservation laws of nucleon number $A$ and charge number $Z$, complete and balance each of the following alpha decay processes:
\begin{enumerate}[leftmargin=1.5em,itemsep=2pt]
  \item Uranium-238 alpha decay: $\ce{^{238}_{92}U ->} \mathbf{?} + \ce{^{4}_{2}He}$
  \item Radium-226 alpha decay: $\ce{^{226}_{88}Ra ->} \mathbf{?} + \ce{^{4}_{2}He}$
  \item Polonium-210 alpha decay: $\ce{^{210}_{84}Po ->} \mathbf{?} + \ce{^{4}_{2}He}$
  \item Americium-241 alpha decay (smoke detectors): $\ce{^{241}_{95}Am ->} \mathbf{?} + \ce{^{4}_{2}He}$
\end{enumerate}
Identify the daughter element in each case using the periodic table.
\end{exercise}

\begin{tabletworkspace}[5.0cm]
\end{tabletworkspace}

\begin{solution}
For alpha decay: $\ce{^{A}_{Z}X -> ^{A-4}_{Z-2}Y + ^{4}_{2}He}$.
\begin{enumerate}[leftmargin=1.5em,itemsep=3pt]
  \item $\ce{^{238}_{92}U}$:
  $A = 238 - 4 = 234,\ Z = 92 - 2 = 90$ (Thorium, $\ce{Th}$):
  \[
  \ce{^{238}_{92}U -> ^{234}_{90}Th + ^{4}_{2}He}.
  \]
  \item $\ce{^{226}_{88}Ra}$:
  $A = 226 - 4 = 222,\ Z = 88 - 2 = 86$ (Radon, $\ce{Rn}$):
  \[
  \ce{^{226}_{88}Ra -> ^{222}_{86}Rn + ^{4}_{2}He}.
  \]
  \item $\ce{^{210}_{84}Po}$:
  $A = 210 - 4 = 206,\ Z = 84 - 2 = 82$ (Lead, $\ce{Pb}$):
  \[
  \ce{^{210}_{84}Po -> ^{206}_{82}Pb + ^{4}_{2}He}.
  \]
  \item $\ce{^{241}_{95}Am}$:
  $A = 241 - 4 = 237,\ Z = 95 - 2 = 93$ (Neptunium, $\ce{Np}$):
  \[
  \ce{^{241}_{95}Am -> ^{237}_{93}Np + ^{4}_{2}He}.
  \]
\end{enumerate}
\end{solution}

% ===========================================================================
% EXERCISE 3.3
% ===========================================================================
\begin{exercise}[$\bigstar$ Foundational]{Balancing Beta Minus ($\beta^-$) and Beta Plus ($\beta^+$) Decays}
Write the complete, balanced nuclear reaction for each decay, including the emitted electron or positron and the corresponding antineutrino ($\bar{\nu}_e$) or neutrino ($\nu_e$):
\begin{enumerate}[leftmargin=1.5em,itemsep=2pt]
  \item Beta minus decay of Carbon-14 ($\ce{^{14}_{6}C}$).
  \item Beta minus decay of Iodine-131 ($\ce{^{131}_{53}I}$).
  \item Beta minus decay of Cobalt-60 ($\ce{^{60}_{27}Co}$).
  \item Beta plus decay (positron emission) of Carbon-11 ($\ce{^{11}_{6}C}$).
  \item Beta plus decay of Fluorine-18 ($\ce{^{18}_{9}F}$) used in PET scanning.
\end{enumerate}
\end{exercise}

\begin{tabletworkspace}[5.0cm]
\end{tabletworkspace}

\begin{solution}
\textbf{Beta Minus Rules:} $\ce{^{A}_{Z}X -> ^{A}_{Z+1}Y + ^{0}_{-1}e + \bar{\nu}_e}$ (with neutron conversion $n \to p + e^- + \bar{\nu}_e$).
\begin{enumerate}[leftmargin=1.5em,itemsep=2pt]
  \item Carbon-14:
  \[
  \ce{^{14}_{6}C -> ^{14}_{7}N + ^{0}_{-1}e + \bar{\nu}_e}.
  \]
  \item Iodine-131:
  \[
  \ce{^{131}_{53}I -> ^{131}_{54}Xe + ^{0}_{-1}e + \bar{\nu}_e}.
  \]
  \item Cobalt-60:
  \[
  \ce{^{60}_{27}Co -> ^{60}_{28}Ni + ^{0}_{-1}e + \bar{\nu}_e}.
  \]
\end{enumerate}

\textbf{Beta Plus Rules:} $\ce{^{A}_{Z}X -> ^{A}_{Z-1}Y + ^{0}_{+1}e + \nu_e}$ (with proton conversion $p \to n + e^+ + \nu_e$).
\begin{enumerate}[leftmargin=1.5em,itemsep=2pt]
  \setcounter{enumi}{3}
  \item Carbon-11:
  \[
  \ce{^{11}_{6}C -> ^{11}_{5}B + ^{0}_{+1}e + \nu_e}.
  \]
  \item Fluorine-18:
  \[
  \ce{^{18}_{9}F -> ^{18}_{8}O + ^{0}_{+1}e + \nu_e}.
  \]
\end{enumerate}
\end{solution}

% ===========================================================================
% EXERCISE 3.4
% ===========================================================================
\begin{exercise}[$\bigstar$ Foundational]{Electron Capture and Gamma De-excitation}
Certain proton-rich nuclides decay competitively via orbital electron capture (EC) rather than positron emission:
\begin{enumerate}[leftmargin=1.5em,itemsep=2pt]
  \item In rocks, Potassium-40 undergoes electron capture by capturing an inner K-shell electron:
  \[
  \ce{^{40}_{19}K + ^{0}_{-1}e ->} \mathbf{?} + \nu_e.
  \]
  Identify the daughter nuclide.
  \item After electron capture, the atomic electron cloud undergoes reorganization. What characteristic electromagnetic radiation is subsequently emitted?
  \item The nuclear isomer Technetium-99m ($\ce{^{99m}_{43}Tc}$) decays by isomeric gamma transition. Write the decay equation and explain why its mass number $A$ and atomic number $Z$ remain unchanged.
\end{enumerate}
\end{exercise}

\begin{tabletworkspace}[5.0cm]
\end{tabletworkspace}

\begin{solution}
\textbf{1. Electron Capture of Potassium-40:}
$A = 40 + 0 = 40,\ Z = 19 + (-1) = 18$ (Argon, $\ce{Ar}$):
\[
\ce{^{40}_{19}K + ^{0}_{-1}e -> ^{40}_{18}Ar + \nu_e}.
\]

\textbf{2. Characteristic Atomic Radiation:}
When an inner K-shell electron is captured by the nucleus, a vacancy is created in the innermost shell. An outer electron (from the L or M shell) drops down into the vacancy, releasing the energy difference as a \textbf{characteristic X-ray photon} ($\ce{K_\alpha}$ or $\ce{K_\beta}$ X-rays) or through the emission of an Auger electron.

\textbf{3. Isomeric Transition of Technetium-99m:}
\[
\ce{^{99m}_{43}Tc -> ^{99}_{43}Tc + \gamma} \quad (E_\gamma = 140.5\ \mathrm{keV}).
\]
The "m" denotes a metastable nuclear excited state. A gamma decay is a pure electromagnetic transition where the nucleus sheds excess excitation energy without altering the number of protons ($Z=43$) or neutrons ($N=56$).
\end{solution}

% ===========================================================================
% EXERCISE 3.5
% ===========================================================================
\begin{exercise}[$\bigstar$ Foundational]{Historic Nuclear Transmutation and Discovery Reactions}
Write the balanced nuclear equations for three monumental discoveries in modern physics:
\begin{enumerate}[leftmargin=1.5em,itemsep=2pt]
  \item \textbf{Discovery of the Proton (Rutherford, 1919):} Nitrogen-14 bombarded with alpha particles produces an unknown nuclide and a free proton ($\ce{^{1}_{1}p}$ or $\ce{^{1}_{1}H}$).
  \item \textbf{Discovery of the Neutron (Chadwick, 1932):} Beryllium-9 bombarded with alpha particles produces Carbon-12 and a neutral particle.
  \item \textbf{First Artificial Radioactivity (Irène \& Frédéric Joliot-Curie, 1934):} Aluminum-27 bombarded with alpha particles ejects a neutron to form a radioactive phosphorus isotope ($\ce{^{30}_{15}P}$), which subsequently decays by positron emission. Write both reactions.
\end{enumerate}
\end{exercise}

\begin{tabletworkspace}[5.0cm]
\end{tabletworkspace}

\begin{solution}
\textbf{1. Rutherford's Proton Discovery (1919):}
Reactants: $A = 14 + 4 = 18,\ Z = 7 + 2 = 9$.
Known product: $\ce{^{1}_{1}p}$.
Daughter nuclide: $A = 18 - 1 = 17,\ Z = 9 - 1 = 8$ (Oxygen, $\ce{O}$):
\[
\ce{^{14}_{7}N + ^{4}_{2}He -> ^{17}_{8}O + ^{1}_{1}p} \quad (\text{Notation: } \ce{^{14}N(\alpha, p)^{17}O}).
\]

\textbf{2. Chadwick's Neutron Discovery (1932):}
Reactants: $A = 9 + 4 = 13,\ Z = 4 + 2 = 6$.
Known product: $\ce{^{12}_{6}C}$.
Neutral particle: $A = 13 - 12 = 1,\ Z = 6 - 6 = 0$ (Neutron, $\ce{^{1}_{0}n}$):
\[
\ce{^{9}_{4}Be + ^{4}_{2}He -> ^{12}_{6}C + ^{1}_{0}n} \quad (\text{Notation: } \ce{^{9}Be(\alpha, n)^{12}C}).
\]

\textbf{3. Joliot-Curie Artificial Radioactivity (1934):}
\begin{itemize}[leftmargin=1.5em,itemsep=2pt]
  \item Bombardment synthesis:
  Reactants: $A = 27 + 4 = 31,\ Z = 13 + 2 = 15$. Ejected particle: $\ce{^{1}_{0}n}$.
  \[
  \ce{^{27}_{13}Al + ^{4}_{2}He -> ^{30}_{15}P + ^{1}_{0}n}.
  \]
  \item Radioactive decay of synthesized phosphorus-30 ($t_{1/2} = 2.50\ \mathrm{min}$):
  \[
  \ce{^{30}_{15}P -> ^{30}_{14}Si + ^{0}_{+1}e + \nu_e}.
  \]
\end{itemize}
\end{solution}

% ===========================================================================
% EXERCISE 3.6
% ===========================================================================
\begin{exercise}[$\bigstar$ Foundational]{Decay Constant, Half-Life, and Mean Lifetime}
The decay kinetics of a radioactive nuclide are governed by the differential equation $\dd N / \dd t = -\lambda N$.
\begin{enumerate}[leftmargin=1.5em,itemsep=2pt]
  \item Integrate this differential equation analytically with initial condition $N(0) = N_0$ to derive $N(t) = N_0 e^{-\lambda t}$.
  \item Derive the exact relationship between the half-life $t_{1/2}$ and the decay constant $\lambda$.
  \item Radon-222 ($\ce{^{222}_{86}Rn}$) has a half-life of $t_{1/2} = 3.823\ \mathrm{days}$. Calculate its decay constant $\lambda$ in $\mathrm{day^{-1}}$ and in SI units ($\mathrm{s^{-1}}$).
  \item Calculate the mean lifetime $\tau = 1/\lambda$ of a Radon-222 nucleus.
\end{enumerate}
\end{exercise}

\begin{tabletworkspace}[5.0cm]
\end{tabletworkspace}

\begin{solution}
\textbf{1. Analytical Integration:}
Separate variables:
\[
\frac{\dd N}{N} = -\lambda \dd t.
\]
Integrate from $t = 0$ ($N = N_0$) to time $t$ ($N = N(t)$):
\[
\int_{N_0}^{N(t)} \frac{\dd N}{N} = -\lambda \int_0^t \dd t \implies \ln \left( \frac{N(t)}{N_0} \right) = -\lambda t.
\]
Exponentiating both sides:
\[
\frac{N(t)}{N_0} = e^{-\lambda t} \implies \mathbf{N(t) = N_0 e^{-\lambda t}}. \quad \text{Q.E.D.}
\]

\textbf{2. Derivation of Half-Life Formula:}
By definition of the half-life $t_{1/2}$, at $t = t_{1/2}$, exactly half the nuclei remain: $N(t_{1/2}) = N_0 / 2$.
\[
\frac{N_0}{2} = N_0 e^{-\lambda t_{1/2}} \implies e^{\lambda t_{1/2}} = 2 \implies \lambda t_{1/2} = \ln 2.
\]
\[
\mathbf{t_{1/2} = \frac{\ln 2}{\lambda} \approx \frac{0.693147}{\lambda}}, \qquad \mathbf{\lambda = \frac{\ln 2}{t_{1/2}}}.
\]

\textbf{3. Calculations for Radon-222:}
\[
\lambda = \frac{\ln 2}{3.823\ \mathrm{days}} = \mathbf{0.1813\ \mathrm{day^{-1}}}.
\]
Converting days to seconds ($1\ \mathrm{day} = 86\,400\ \mathrm{s}$):
\[
t_{1/2} = 3.823 \times 86\,400\ \mathrm{s} = 330\,307.2\ \mathrm{s}.
\]
\[
\lambda = \frac{\ln 2}{330\,307.2\ \mathrm{s}} = \mathbf{2.0985 \times 10^{-6}\ \mathrm{s^{-1}}}.
\]

\textbf{4. Mean Lifetime $\tau$:}
\[
\tau = \frac{1}{\lambda} = \frac{1}{0.1813\ \mathrm{day^{-1}}} = \mathbf{5.515\ \mathrm{days}} \approx \mathbf{4.765 \times 10^5\ \mathrm{s}}.
\]
Notice that $\tau = t_{1/2} / \ln 2 \approx 1.443 \, t_{1/2}$.
\end{solution}

% ===========================================================================
% EXERCISE 3.7
% ===========================================================================
\begin{exercise}[$\bigstar$ Foundational]{Absolute Activity: The Becquerel and the Curie}
The historical Curie unit ($1\ \mathrm{Ci}$) was originally defined by Marie Curie as the activity of exactly $1.000\ \mathrm{g}$ of pure Radium-226 ($\ce{^{226}_{88}Ra}$).
\begin{itemize}[itemsep=1pt]
  \item Molar mass of Radium-226: $M = 226.025\ \mathrm{g\cdot mol^{-1}}$
  \item Half-life: $t_{1/2} = 1600\ \mathrm{years}$ ($1\ \mathrm{yr} = 365.25\ \mathrm{days} = 3.15576 \times 10^7\ \mathrm{s}$)
  \item $N_A = 6.022\,141 \times 10^{23}\ \mathrm{mol^{-1}}$
\end{itemize}
\begin{enumerate}[leftmargin=1.5em,itemsep=2pt]
  \item Calculate the decay constant $\lambda$ of Radium-226 in $\mathrm{s^{-1}}$.
  \item Determine the number of radium nuclei $N$ in $1.000\ \mathrm{g}$ of $\ce{^{226}Ra}$.
  \item Compute the absolute activity $A = \lambda N$ in Becquerels ($\mathrm{Bq}$).
  \item Compare your result to the modern exact definition: $1\ \mathrm{Ci} \equiv 3.700 \times 10^{10}\ \mathrm{Bq}$.
\end{enumerate}
\end{exercise}

\begin{tabletworkspace}[5.0cm]
\end{tabletworkspace}

\begin{solution}
\textbf{1. Decay Constant $\lambda$ in $\mathrm{s^{-1}}$:}
Half-life in seconds:
\[
t_{1/2} = 1600 \times 3.15576 \times 10^7\ \mathrm{s} = 5.0492 \times 10^{10}\ \mathrm{s}.
\]
\[
\lambda = \frac{\ln 2}{t_{1/2}} = \frac{0.693147}{5.0492 \times 10^{10}\ \mathrm{s}} = \mathbf{1.3728 \times 10^{-11}\ \mathrm{s^{-1}}}.
\]

\textbf{2. Number of Radium Nuclei in $1.000\ \mathrm{g}$:}
\[
N = \frac{m}{M} \cdot N_A = \frac{1.000\ \mathrm{g}}{226.025\ \mathrm{g\cdot mol^{-1}}} \times 6.022\,141 \times 10^{23}\ \mathrm{mol^{-1}} = \mathbf{2.6644 \times 10^{21}\ \text{nuclei}}.
\]

\textbf{3. Absolute Activity $A$ in Becquerels:}
\[
A = \lambda N = (1.3728 \times 10^{-11}\ \mathrm{s^{-1}}) \times (2.6644 \times 10^{21}) = \mathbf{3.658 \times 10^{10}\ \mathrm{Bq}}.
\]

\textbf{4. Comparison with Standard Definition:}
The calculated activity of $1\ \mathrm{g}$ of $^{226}\ce{Ra}$ ($3.66 \times 10^{10}\ \mathrm{Bq}$) differs by only $1.1\%$ from the standardized definition of the Curie ($3.700 \times 10^{10}\ \mathrm{Bq}$), which was established internationally in 1953 as an exact rounded reference benchmark.
\end{solution}

% ===========================================================================
% EXERCISE 3.8
% ===========================================================================
\begin{exercise}[$\bigstar\bigstar$ Intermediate]{Radioactive Mass Decay of Iodine-131}
Iodine-131 ($\ce{^{131}_{53}I}$) is widely used in hospital nuclear medicine for the therapeutic destruction of hyperactive thyroid tissue. It decays by $\beta^-$ emission with a half-life of $t_{1/2} = 8.020\ \mathrm{days}$. A hospital pharmacy receives an initial sample containing $m_0 = 20.00\ \mathrm{mg}$ of pure $\ce{^{131}I}$.
\begin{enumerate}[leftmargin=1.5em,itemsep=2pt]
  \item Calculate the remaining mass of $\ce{^{131}I}$ after:
  \begin{enumerate}[leftmargin=1.2em,itemsep=1pt]
    \item Exactly two half-lives ($t = 16.04\ \mathrm{days}$).
    \item Exactly three half-lives ($t = 24.06\ \mathrm{days}$).
    \item Exactly $30.00\ \mathrm{days}$.
  \end{enumerate}
  \item What percentage of the initial radioactivity has decayed after $30.00\ \mathrm{days}$?
  \item How many days must elapse until only $0.100\%$ ($1/1000$) of the initial iodine remains?
\end{enumerate}
\end{exercise}

\begin{tabletworkspace}[5.0cm]
\end{tabletworkspace}

\begin{solution}
\textbf{1. Remaining Mass:}
\begin{enumerate}[leftmargin=1.4em,itemsep=2pt]
  \item After $k = 2$ half-lives:
  \[
  m(16.04\ \mathrm{d}) = \frac{m_0}{2^2} = \frac{20.00\ \mathrm{mg}}{4} = \mathbf{5.000\ \mathrm{mg}}.
  \]
  \item After $k = 3$ half-lives:
  \[
  m(24.06\ \mathrm{d}) = \frac{m_0}{2^3} = \frac{20.00\ \mathrm{mg}}{8} = \mathbf{2.500\ \mathrm{mg}}.
  \]
  \item After $t = 30.00\ \mathrm{days}$:
  First find $\lambda$: $\lambda = \frac{\ln 2}{8.020\ \mathrm{d}} = 0.086427\ \mathrm{day^{-1}}$.
  \[
  m(30.0\ \mathrm{d}) = m_0 e^{-\lambda t} = 20.00 \times e^{-0.086427 \times 30.00} = 20.00 \times e^{-2.5928} = \mathbf{1.496\ \mathrm{mg}}.
  \]
\end{enumerate}

\textbf{2. Percentage Decayed after $30.00\ \mathrm{days}$:}
Fraction remaining: $\frac{1.496}{20.00} = 0.0748 \implies 7.48\%$.
\[
\text{Percentage Decayed} = 100\% - 7.48\% = \mathbf{92.52\%}.
\]

\textbf{3. Time for $0.100\%$ Remaining ($N(t)/N_0 = 10^{-3}$):}
\[
\frac{N(t)}{N_0} = e^{-\lambda t} = 0.00100 \implies -\lambda t = \ln(0.00100) = -6.907755.
\]
\[
t = \frac{6.907755}{0.086427\ \mathrm{day^{-1}}} = \mathbf{79.93\ \mathrm{days}} \approx \mathbf{80\ \mathrm{days}}.
\]
(Rule of thumb: 10 half-lives reduces activity by $2^{10} = 1024 \approx 1000$; here $10 \times 8.02 = 80.2\ \mathrm{days}$).
\end{solution}

% ===========================================================================
% EXERCISE 3.9
% ===========================================================================
\begin{exercise}[$\bigstar\bigstar$ Intermediate]{Determination of Radioactive Half-Life from Activity Data}
A laboratory technician measures the count rate (activity in counts per minute, $\mathrm{cpm}$) of a newly isolated beta-emitting radioisotope using a Geiger-Müller counter. The recorded data are:
\begin{center}
\begin{tabular}{lcccccc}
\toprule
Time $t$ (hours) & $0$ & $3.0$ & $6.0$ & $9.0$ & $12.0$ & $18.0$ \\
\midrule
Activity $A(t)$ ($\mathrm{cpm}$) & $8400$ & $5940$ & $4200$ & $2970$ & $2100$ & $1050$ \\
\bottomrule
\end{tabular}
\end{center}
\begin{enumerate}[leftmargin=1.5em,itemsep=2pt]
  \item By inspecting the activity values at $t = 0\ \mathrm{h}$, $t = 6.0\ \mathrm{h}$, $t = 12.0\ \mathrm{h}$, and $t = 18.0\ \mathrm{h}$, deduce the half-life $t_{1/2}$ directly.
  \item Calculate the decay constant $\lambda$ in $\mathrm{h^{-1}}$ and in $\mathrm{s^{-1}}$.
  \item Predict the expected activity of the sample at $t = 24.0\ \mathrm{hours}$ and at $t = 48.0\ \mathrm{hours}$.
\end{enumerate}
\end{exercise}

\begin{tabletworkspace}[5.0cm]
\end{tabletworkspace}

\begin{solution}
\textbf{1. Direct Deduction of Half-Life:}
Observe the sequence of recorded activities:
\begin{itemize}[itemsep=1pt]
  \item At $t = 0\ \mathrm{h}$: $A_0 = 8400\ \mathrm{cpm}$.
  \item At $t = 6.0\ \mathrm{h}$: $A(6) = 4200\ \mathrm{cpm} = A_0 / 2$.
  \item At $t = 12.0\ \mathrm{h}$: $A(12) = 2100\ \mathrm{cpm} = A_0 / 4$.
  \item At $t = 18.0\ \mathrm{h}$: $A(18) = 1050\ \mathrm{cpm} = A_0 / 8$.
\end{itemize}
In every interval of $6.0\ \mathrm{hours}$, the activity is halved precisely. Therefore:
\[
\mathbf{t_{1/2} = 6.00\ \mathrm{hours}}.
\]
(This corresponds to the well-known medical radioisotope Technetium-99m!).

\textbf{2. Decay Constant $\lambda$:}
\[
\lambda = \frac{\ln 2}{6.00\ \mathrm{h}} = \mathbf{0.1155\ \mathrm{h^{-1}}}.
\]
In SI units ($1\ \mathrm{h} = 3600\ \mathrm{s}$):
\[
\lambda = \frac{0.115525\ \mathrm{h^{-1}}}{3600\ \mathrm{s\cdot h^{-1}}} = \mathbf{3.209 \times 10^{-5}\ \mathrm{s^{-1}}}.
\]

\textbf{3. Predicted Activities:}
\begin{itemize}[itemsep=1pt]
  \item At $t = 24.0\ \mathrm{h}$ ($4$ half-lives):
  \[
  A(24) = \frac{A_0}{2^4} = \frac{8400}{16} = \mathbf{525\ \mathrm{cpm}}.
  \]
  \item At $t = 48.0\ \mathrm{h}$ ($8$ half-lives):
  \[
  A(48) = \frac{A_0}{2^8} = \frac{8400}{256} = \mathbf{32.8\ \mathrm{cpm}}.
  \]
\end{itemize}
\end{solution}

% ===========================================================================
% EXERCISE 3.10
% ===========================================================================
\begin{exercise}[$\bigstar\bigstar$ Intermediate]{Specific Activity of Pure Phosphorus-32}
Phosphorus-32 ($\ce{^{32}_{15}P}$, $M = 31.974\ \mathrm{g\cdot mol^{-1}}$) is a pure $\beta^-$ emitter with $t_{1/2} = 14.26\ \mathrm{days}$, used in molecular biology for radiolabeling nucleic acids.
\begin{enumerate}[leftmargin=1.5em,itemsep=2pt]
  \item Define the \textbf{specific activity} ($a$) of a radioactive source.
  \item Derive the algebraic formula for the specific activity $a$ of a carrier-free pure radionuclide in terms of $N_A, M,$ and $t_{1/2}$.
  \item Calculate the specific activity of carrier-free pure $\ce{^{32}P}$ in:
  \begin{enumerate}[leftmargin=1.2em,itemsep=1pt]
    \item Becquerels per gram ($\mathrm{Bq\cdot g^{-1}}$).
    \item Curies per gram ($\mathrm{Ci\cdot g^{-1}}$).
  \end{enumerate}
\end{enumerate}
\end{exercise}

\begin{tabletworkspace}[5.0cm]
\end{tabletworkspace}

\begin{solution}
\textbf{1. Definition of Specific Activity:}
The specific activity $a$ is the activity per unit mass of the substance:
\[
a = \frac{A}{m} \quad [\mathrm{Bq\cdot g^{-1}} \text{ or } \mathrm{Ci\cdot g^{-1}}].
\]

\textbf{2. Algebraic Derivation:}
For a sample of mass $m$ of a pure radionuclide with molar mass $M$:
\[
N = \frac{m}{M} \cdot N_A.
\]
The activity is $A = \lambda N = \frac{\ln 2}{t_{1/2}} \cdot \frac{m}{M} \cdot N_A$.
Dividing by $m$:
\[
\mathbf{a = \frac{\ln 2 \cdot N_A}{M \cdot t_{1/2}}}. \quad \text{Q.E.D.}
\]

\textbf{3. Numerical Calculation for Phosphorus-32:}
Half-life in seconds:
\[
t_{1/2} = 14.26 \times 86\,400\ \mathrm{s} = 1.232\,064 \times 10^6\ \mathrm{s}.
\]
\begin{enumerate}[leftmargin=1.4em,itemsep=2pt]
  \item Specific activity in $\mathrm{Bq\cdot g^{-1}}$:
  \begin{align*}
  a &= \frac{0.693\,147 \times 6.022\,141 \times 10^{23}\ \mathrm{mol^{-1}}}{31.974\ \mathrm{g\cdot mol^{-1}} \times 1.232\,064 \times 10^6\ \mathrm{s}} \\
    &= \frac{4.174\,23 \times 10^{23}}{3.9394 \times 10^7} = \mathbf{1.0596 \times 10^{16}\ \mathrm{Bq\cdot g^{-1}}}.
  \end{align*}
  \item Specific activity in $\mathrm{Ci\cdot g^{-1}}$:
  \[
  a = \frac{1.0596 \times 10^{16}\ \mathrm{Bq\cdot g^{-1}}}{3.700 \times 10^{10}\ \mathrm{Bq\cdot Ci^{-1}}} = \mathbf{286\,380\ \mathrm{Ci\cdot g^{-1}}} \approx \mathbf{2.86 \times 10^5\ \mathrm{Ci\cdot g^{-1}}}.
  \]
  A single milligram of carrier-free $\ce{^{32}P}$ carries an astonishing activity of $286\ \mathrm{Ci}$!
\end{enumerate}
\end{solution}

% ===========================================================================
% EXERCISE 3.11
% ===========================================================================
\begin{exercise}[$\bigstar\bigstar$ Intermediate]{Nuclear Medicine Radiopharmaceutical Dosage Protocol}
A nuclear medicine department orders Technetium-99m ($\ce{^{99m}_{43}Tc}$, $t_{1/2} = 6.010\ \mathrm{hours}$) from a regional radiopharmacy. The radiopharmaceutical is calibrated to have an activity of $A_0 = 740.0\ \mathrm{MBq}$ ($20.00\ \mathrm{mCi}$) upon delivery at the hospital at $08:00\ \mathrm{AM}$.
\begin{enumerate}[leftmargin=1.5em,itemsep=2pt]
  \item A patient is scheduled for a cardiac perfusion SPECT scan at $14:30$ ($2:30\ \mathrm{PM}$) on the same day. Calculate the activity $A$ remaining in the syringe at the time of injection.
  \item If the scan protocol prescribes a required injected dose of exactly $400.0\ \mathrm{MBq}$, at what exact time of day should the patient have been injected with this syringe?
  \item If unused, how much activity remains in the syringe after $24.0\ \mathrm{hours}$ (at $08:00\ \mathrm{AM}$ the next day)?
\end{enumerate}
\end{exercise}

\begin{tabletworkspace}[5.0cm]
\end{tabletworkspace}

\begin{solution}
\textbf{1. Activity at $14:30$ (Elapsed time $\Delta t = 6.50\ \mathrm{hours}$):}
Decay constant:
\[
\lambda = \frac{\ln 2}{6.010\ \mathrm{h}} = 0.115332\ \mathrm{h^{-1}}.
\]
Remaining activity:
\[
A(6.50\ \mathrm{h}) = A_0 e^{-\lambda t} = 740.0\ \mathrm{MBq} \times e^{-0.115332 \times 6.50} = 740.0 \times e^{-0.74966} = \mathbf{349.5\ \mathrm{MBq}}.
\]

\textbf{2. Injection Time for Exactly $400.0\ \mathrm{MBq}$:}
\[
\frac{A(t)}{A_0} = \frac{400.0}{740.0} = 0.54054 \implies e^{-\lambda t} = 0.54054
\]
\[
-\lambda t = \ln(0.54054) = -0.61519 \implies t = \frac{0.61519}{0.115332\ \mathrm{h^{-1}}} = 5.334\ \mathrm{hours}.
\]
Convert decimal hours to hours and minutes:
\[
0.334\ \mathrm{h} \times 60\ \mathrm{min/h} = 20.04\ \mathrm{min} \approx 20\ \mathrm{minutes}.
\]
Adding $5\ \mathrm{hours}\ 20\ \mathrm{minutes}$ to $08:00\ \mathrm{AM}$:
\[
\text{Injection Time} = \mathbf{13:20\ (1:20\ \mathrm{PM})}.
\]

\textbf{3. Residual Activity After 24.0 Hours ($k = 24.0 / 6.010 = 3.993 \approx 4$ half-lives):}
\[
A(24\ \mathrm{h}) = 740.0 \times e^{-0.115332 \times 24.0} = 740.0 \times e^{-2.7680} = \mathbf{46.33\ \mathrm{MBq}} \approx \mathbf{1.25\ \mathrm{mCi}}.
\]
\end{solution}

% ===========================================================================
% EXERCISE 3.12
% ===========================================================================
\begin{exercise}[$\bigstar\bigstar$ Intermediate]{Counting Particle Emissions in Natural Decay Series}
The natural radioactive decay series cascade through multiple intermediate radionuclides before terminating at a stable lead isotope ($\ce{Pb}$, $Z = 82$).
\begin{enumerate}[leftmargin=1.5em,itemsep=2pt]
  \item \textbf{Thorium Series ($4n$):} Thorium-232 ($\ce{^{232}_{90}Th}$) decays through a series of $\alpha$ and $\beta^-$ emissions to stable Lead-208 ($\ce{^{208}_{82}Pb}$).
  Determine the exact number of alpha particles ($N_\alpha$) and beta particles ($N_\beta$) emitted per transformed parent atom.
  \item \textbf{Actinium Series ($4n+3$):} Uranium-235 ($\ce{^{235}_{92}U}$) decays to stable Lead-207 ($\ce{^{207}_{82}Pb}$).
  Determine $N_\alpha$ and $N_\beta$.
  \item State the general algebraic formula for $N_\alpha$ and $N_\beta$ in any decay chain transitioning from $\ce{^{A_1}_{Z_1}X}$ to $\ce{^{A_2}_{Z_2}Y}$.
\end{enumerate}
\end{exercise}

\begin{tabletworkspace}[5.0cm]
\end{tabletworkspace}

\begin{solution}
\textbf{3. General Algebraic Formula:}
Let the overall chain transformation be:
\[
\ce{^{A_1}_{Z_1}X -> ^{A_2}_{Z_2}Y + N_\alpha \, ^{4}_{2}He + N_\beta \, ^{0}_{-1}e}.
\]
Conservation of mass number $A$:
\[
A_1 = A_2 + 4 N_\alpha + 0 \implies \mathbf{N_\alpha = \frac{A_1 - A_2}{4}}.
\]
Conservation of charge number $Z$:
\[
Z_1 = Z_2 + 2 N_\alpha - N_\beta \implies \mathbf{N_\beta = 2 N_\alpha - (Z_1 - Z_2)}.
\]

\textbf{1. Thorium Series ($\ce{^{232}_{90}Th -> ^{208}_{82}Pb}$):}
\begin{align*}
N_\alpha &= \frac{232 - 208}{4} = \frac{24}{4} = \mathbf{6\ \alpha\text{-particles}}. \\
N_\beta &= 2(6) - (90 - 82) = 12 - 8 = \mathbf{4\ \beta^-\text{-particles}}.
\end{align*}
Overall balanced equation:
\[
\ce{^{232}_{90}Th -> ^{208}_{82}Pb + 6\,^{4}_{2}He + 4\,^{0}_{-1}e + 4\,\bar{\nu}_e}.
\]

\textbf{2. Actinium Series ($\ce{^{235}_{92}U -> ^{207}_{82}Pb}$):}
\begin{align*}
N_\alpha &= \frac{235 - 207}{4} = \frac{28}{4} = \mathbf{7\ \alpha\text{-particles}}. \\
N_\beta &= 2(7) - (92 - 82) = 14 - 10 = \mathbf{4\ \beta^-\text{-particles}}.
\end{align*}
Overall balanced equation:
\[
\ce{^{235}_{92}U -> ^{207}_{82}Pb + 7\,^{4}_{2}He + 4\,^{0}_{-1}e + 4\,\bar{\nu}_e}.
\]
\end{solution}

% ===========================================================================
% EXERCISE 3.13
% ===========================================================================
\begin{exercise}[$\bigstar\bigstar$ Intermediate]{Helium Accumulation from Alpha Emitters}
Alpha particles are energetic helium nuclei that capture two stray electrons from surrounding matter to become neutral gaseous Helium-4 ($\ce{^{4}_{2}He}$).
A sealed ampoule contains $m = 1.000\ \mathrm{g}$ of pure Radium-226 chloride ($\ce{RaCl2}$, $M = 296.93\ \mathrm{g\cdot mol^{-1}}$). The radium decays by alpha emission with a half-life of $t_{1/2} = 1600\ \mathrm{years}$.
\begin{enumerate}[leftmargin=1.5em,itemsep=2pt]
  \item Compute the number of $\ce{^{226}Ra}$ nuclei initially present in the ampoule.
  \item Assuming the ampoule remains sealed for exactly $t = 10.0\ \mathrm{years}$, calculate the number of radium nuclei that decay.
  \item Assuming each decaying radium nucleus yields one helium atom, determine the volume of helium gas accumulated inside the ampoule measured at standard temperature and pressure ($\mathrm{STP}: V_0 = 22.414\ \mathrm{L\cdot mol^{-1}}$).
\end{enumerate}
\end{exercise}

\begin{tabletworkspace}[5.0cm]
\end{tabletworkspace}

\begin{solution}
\textbf{1. Number of Initial Radium Nuclei ($N_0$):}
\[
n(\ce{RaCl2}) = \frac{1.000\ \mathrm{g}}{296.93\ \mathrm{g\cdot mol^{-1}}} = 3.3678 \times 10^{-3}\ \mathrm{mol}.
\]
Since each formula unit contains one radium atom:
\[
N_0 = 3.3678 \times 10^{-3}\ \mathrm{mol} \times 6.022\,141 \times 10^{23}\ \mathrm{mol^{-1}} = \mathbf{2.0281 \times 10^{21}\ \text{nuclei}}.
\]

\textbf{2. Number of Decays in $10.0\ \mathrm{years}$:}
Decay constant in $\mathrm{yr^{-1}}$:
\[
\lambda = \frac{\ln 2}{1600\ \mathrm{yr}} = 4.332 \times 10^{-4}\ \mathrm{yr^{-1}}.
\]
Since $\lambda t = 4.332 \times 10^{-4} \times 10.0 = 4.332 \times 10^{-3} \ll 1$, we can use the first-order Taylor expansion $1 - e^{-\lambda t} \approx \lambda t$:
\[
\Delta N = N_0(1 - e^{-\lambda t}) \approx N_0 \cdot \lambda \cdot t.
\]
\[
\Delta N = (2.0281 \times 10^{21}) \times (4.332 \times 10^{-3}) = \mathbf{8.786 \times 10^{18}\ \text{decays}}.
\]

\textbf{3. Volume of Helium Gas at STP:}
Number of moles of helium gas formed:
\[
n(\ce{He}) = \frac{\Delta N}{N_A} = \frac{8.786 \times 10^{18}}{6.022\,141 \times 10^{23}} = 1.4589 \times 10^{-5}\ \mathrm{mol}.
\]
Volume at STP:
\[
V(\ce{He}) = n \cdot V_0 = 1.4589 \times 10^{-5}\ \mathrm{mol} \times 22.414\ \mathrm{L\cdot mol^{-1}} = 3.270 \times 10^{-4}\ \mathrm{L} = \mathbf{0.327\ \mathrm{mL}}.
\]
Over a single decade, one gram of radium chloride generates nearly a third of a milliliter of pure helium gas!
\end{solution}

% ===========================================================================
% EXERCISE 3.14
% ===========================================================================
\begin{exercise}[$\bigstar\bigstar$ Intermediate]{Secular Equilibrium in the Uranium Decay Series}
In natural uranium minerals that have remained undisturbed for geological epochs, parent Uranium-238 ($\ce{^{238}U}$, $t_{1/2,1} = 4.468 \times 10^9\ \mathrm{yr}$) exists in \textbf{secular equilibrium} with its immediate daughter Thorium-234 ($\ce{^{234}Th}$, $t_{1/2,2} = 24.10\ \mathrm{days}$).
\begin{enumerate}[leftmargin=1.5em,itemsep=2pt]
  \item State the condition of secular equilibrium and show why the activities of parent and daughter are equal: $A_1 = A_2$.
  \item Derive the equilibrium mass ratio $\dfrac{m_2}{m_1} = \dfrac{t_{1/2,2}}{t_{1/2,1}} \cdot \dfrac{M_2}{M_1}$.
  \item Calculate the mass of Thorium-234 present in secular equilibrium with $m_1 = 1.000\ \mathrm{kg}$ of natural Uranium-238.
\end{enumerate}
\end{exercise}

\begin{tabletworkspace}[5.0cm]
\end{tabletworkspace}

\begin{solution}
\textbf{1. Secular Equilibrium Condition:}
When the parent has a half-life enormously longer than the daughter ($t_{1/2,1} \gg t_{1/2,2}$, here $4.5 \times 10^9\ \mathrm{yr} \gg 24\ \mathrm{days}$), the parent's decay rate is virtually constant. The rate of creation of daughter nuclei ($\lambda_1 N_1$) balances the rate of daughter decay ($\lambda_2 N_2$):
\[
\frac{\dd N_2}{\dd t} = \lambda_1 N_1 - \lambda_2 N_2 = 0 \implies \mathbf{\lambda_1 N_1 = \lambda_2 N_2} \iff \mathbf{A_1 = A_2}.
\]

\textbf{2. Derivation of Equilibrium Mass Ratio:}
From $\lambda_1 N_1 = \lambda_2 N_2$:
\[
\frac{N_2}{N_1} = \frac{\lambda_1}{\lambda_2} = \frac{\ln 2 / t_{1/2,1}}{\ln 2 / t_{1/2,2}} = \frac{t_{1/2,2}}{t_{1/2,1}}.
\]
Substitute $N = \frac{m}{M} N_A$:
\[
\frac{m_2 / M_2}{m_1 / M_1} = \frac{t_{1/2,2}}{t_{1/2,1}} \implies \mathbf{\frac{m_2}{m_1} = \frac{t_{1/2,2}}{t_{1/2,1}} \cdot \frac{M_2}{M_1}}. \quad \text{Q.E.D.}
\]

\textbf{3. Calculation of Equilibrium Mass of $^{234}\ce{Th}$:}
Convert half-lives to identical units (days):
\[
t_{1/2,1} = 4.468 \times 10^9 \times 365.25\ \mathrm{days} = 1.632 \times 10^{12}\ \mathrm{days}.
\]
\[
m_2 = m_1 \times \left( \frac{24.10\ \mathrm{days}}{1.632 \times 10^{12}\ \mathrm{days}} \right) \times \left( \frac{234.04}{238.05} \right) = 1000\ \mathrm{g} \times (1.4767 \times 10^{-11}) \times 0.98315.
\]
\[
m_2 = \mathbf{1.452 \times 10^{-8}\ \mathrm{g}} = \mathbf{14.52\ \mathrm{ng}}.
\]
In an entire kilogram of uranium ore, only $14.5\ \mathrm{nanograms}$ of thorium-234 exists at any given instant, yet its activity matches that of the entire kilogram of uranium!
\end{solution}

% ===========================================================================
% EXERCISE 3.15
% ===========================================================================
\begin{exercise}[$\bigstar\bigstar$ Intermediate]{Radiocarbon Archaeological Dating of Ancient Wood}
Living plant matter maintains a constant dynamic equilibrium with atmospheric Carbon-14, exhibiting a specific activity of $A_0 = 15.30$ disintegrations per minute per gram of total carbon ($\mathrm{dpm\cdot g^{-1}}$). After the organism dies, carbon exchange ceases and $\ce{^{14}C}$ decays with a half-life of $t_{1/2} = 5730\ \mathrm{years}$.
\begin{enumerate}[leftmargin=1.5em,itemsep=2pt]
  \item A wooden beam excavated from the tomb of an Egyptian pharaoh yields a measured specific activity of $A(t) = 9.42\ \mathrm{dpm\cdot g^{-1}}$.
  \item Derive the dating formula expressing age $t$ in terms of $t_{1/2}, A_0,$ and $A(t)$.
  \item Calculate the age of the wooden beam to the nearest decade.
  \item In what calendar era (BCE) was the tree felled?
\end{enumerate}
\end{exercise}

\begin{tabletworkspace}[5.0cm]
\end{tabletworkspace}

\begin{solution}
\textbf{1 \& 2. Derivation of Radiocarbon Dating Formula:}
From the decay equation $A(t) = A_0 e^{-\lambda t}$:
\[
\frac{A(t)}{A_0} = e^{-\lambda t} \implies \ln \left( \frac{A(t)}{A_0} \right) = -\lambda t \implies \ln \left( \frac{A_0}{A(t)} \right) = \lambda t.
\]
Since $\lambda = \dfrac{\ln 2}{t_{1/2}}$:
\[
\mathbf{t = \frac{1}{\lambda} \ln \left( \frac{A_0}{A(t)} \right) = \frac{t_{1/2}}{\ln 2} \ln \left( \frac{A_0}{A(t)} \right)}. \quad \text{Q.E.D.}
\]

\textbf{3. Numerical Calculation of Age:}
Given $A_0 = 15.30\ \mathrm{dpm\cdot g^{-1}}$, $A(t) = 9.42\ \mathrm{dpm\cdot g^{-1}}$, and $t_{1/2} = 5730\ \mathrm{yr}$:
\[
\frac{A_0}{A(t)} = \frac{15.30}{9.42} = 1.6242.
\]
\[
\ln(1.6242) = 0.48501.
\]
\[
t = \frac{5730\ \mathrm{yr}}{0.693147} \times 0.48501 = 8266.6 \times 0.48501 = \mathbf{4009\ \mathrm{years}} \approx \mathbf{4010\ \mathrm{years}}.
\]

\textbf{4. Calendar Era:}
Subtracting current date ($2026$):
\[
\text{Date} = 2026 - 4010 = -1984 \implies \mathbf{\approx 1980\ \mathrm{BCE}} \quad \text{(Middle Kingdom of Ancient Egypt)}.
\]
\end{solution}

% ===========================================================================
% EXERCISE 3.16
% ===========================================================================
\begin{exercise}[$\bigstar\bigstar\bigstar$ Advanced]{Dating with Measurement Uncertainty Propagation}
Charcoal from the prehistoric Chauvet Cave (France) is analyzed by accelerator mass spectrometry (AMS). The remaining $\ce{^{14}C}$ activity is measured to be $A = 1.85 \pm 0.05\%$ of modern reference carbon ($A_0 = 100\%$).
\begin{enumerate}[leftmargin=1.5em,itemsep=2pt]
  \item Calculate the central radiocarbon age $t$ of the prehistoric cave charcoal ($t_{1/2} = 5730\ \mathrm{yr}$).
  \item Using first-order differential error propagation, show that the uncertainty in age is given by $\Delta t = \dfrac{t_{1/2}}{\ln 2} \cdot \dfrac{\Delta A}{A}$.
  \item Calculate the absolute uncertainty $\Delta t$ and express the final dated age in standard scientific format: $t \pm \Delta t\ \mathrm{years\ BP}$ (Before Present).
\end{enumerate}
\end{exercise}

\begin{tabletworkspace}[5.0cm]
\end{tabletworkspace}

\begin{solution}
\textbf{1. Central Age Calculation:}
Activity ratio: $\frac{A}{A_0} = 0.0185 \implies \frac{A_0}{A} = \frac{1}{0.0185} = 54.054$.
\[
t = \frac{5730}{\ln 2} \ln(54.054) = 8266.64 \times 3.9900 = \mathbf{32\,984\ \mathrm{years\ BP}}.
\]

\textbf{2. Differential Error Propagation Derivation:}
Differentiating the dating equation $t = -\frac{1}{\lambda} \ln A + \text{const}$ with respect to $A$:
\[
\frac{\dd t}{\dd A} = -\frac{1}{\lambda A} \implies |\Delta t| = \frac{1}{\lambda} \left| \frac{\Delta A}{A} \right| = \mathbf{\frac{t_{1/2}}{\ln 2} \cdot \frac{\Delta A}{A}}. \quad \text{Q.E.D.}
\]

\textbf{3. Uncertainty Calculation:}
Given $A = 1.85$ and $\Delta A = 0.05$:
\[
\frac{\Delta A}{A} = \frac{0.05}{1.85} = 0.02703 \quad (2.70\%).
\]
\[
\Delta t = 8266.64 \times 0.02703 = \mathbf{223\ \mathrm{years}}.
\]
Final calibrated age:
\[
\mathbf{t = 33\,000 \pm 220\ \mathrm{years\ BP}} \quad \text{(Aurignacian paleolithic era)}.
\]
\end{solution}

% ===========================================================================
% EXERCISE 3.17
% ===========================================================================
\begin{exercise}[$\bigstar\bigstar\bigstar$ Advanced]{Uranium-Lead Geochronology of Ancient Zircons}
Zircon crystals ($\ce{ZrSiO4}$) incorporate uranium atoms into their crystal lattice during crystallization while rigorously rejecting lead ($\ce{Pb^{2+}}$) due to ionic radius mismatch ($r(\ce{Zr^{4+}}) = 84\ \mathrm{pm}$ vs $r(\ce{Pb^{2+}}) = 119\ \mathrm{pm}$). Thus, any $\ce{^{206}Pb}$ present today arose exclusively from the alpha decay cascade of Uranium-238 ($\ce{^{238}U} \to \ce{^{206}Pb}$, $t_{1/2} = 4.468 \times 10^9\ \mathrm{yr}$).
\begin{enumerate}[leftmargin=1.5em,itemsep=2pt]
  \item Derive the geochronological dating equation:
  \[
  t = \frac{1}{\lambda} \ln \left( 1 + \frac{N_{\ce{Pb}}}{N_{\ce{U}}} \right).
  \]
  \item Mass spectrometric analysis of an ancient zircon crystal from the Jack Hills (Western Australia) yields an isotopic atom ratio of:
  \[
  \frac{N(\ce{^{206}Pb})}{N(\ce{^{238}U})} = 0.9635.
  \]
  Calculate the crystallization age $t$ of this zircon crystal in billions of years ($\mathrm{Ga}$).
  \item Compare this age to the accepted scientific age of the Earth ($4.54\ \mathrm{Ga}$).
\end{enumerate}
\end{exercise}

\begin{tabletworkspace}[5.0cm]
\end{tabletworkspace}

\begin{solution}
\textbf{1. Derivation of Uranium-Lead Isochron Formula:}
Let $N_0$ be the initial number of uranium nuclei at crystallization ($t = 0$).
Since no initial lead was incorporated ($N_{\ce{Pb}}(0) = 0$), every radiogenic $\ce{^{206}Pb}$ nucleus formed corresponds to one decayed $\ce{^{238}U}$ nucleus:
\[
N_{\ce{Pb}}(t) = N_0 - N_{\ce{U}}(t).
\]
From exponential decay, $N_0 = N_{\ce{U}}(t) e^{\lambda t}$. Substituting:
\[
N_{\ce{Pb}}(t) = N_{\ce{U}}(t) e^{\lambda t} - N_{\ce{U}}(t) = N_{\ce{U}}(t) \left( e^{\lambda t} - 1 \right).
\]
Dividing by $N_{\ce{U}}(t)$:
\[
\frac{N_{\ce{Pb}}}{N_{\ce{U}}} = e^{\lambda t} - 1 \implies e^{\lambda t} = 1 + \frac{N_{\ce{Pb}}}{N_{\ce{U}}}.
\]
Taking natural logarithms:
\[
\mathbf{t = \frac{1}{\lambda} \ln \left( 1 + \frac{N_{\ce{Pb}}}{N_{\ce{U}}} \right)}. \quad \text{Q.E.D.}
\]

\textbf{2. Age Calculation of Jack Hills Zircon:}
Decay constant of $\ce{^{238}U}$:
\[
\lambda = \frac{\ln 2}{4.468 \times 10^9\ \mathrm{yr}} = 1.55136 \times 10^{-10}\ \mathrm{yr^{-1}}.
\]
Given ratio $N_{\ce{Pb}} / N_{\ce{U}} = 0.9635$:
\[
1 + \frac{N_{\ce{Pb}}}{N_{\ce{U}}} = 1 + 0.9635 = 1.9635.
\]
\[
\ln(1.9635) = 0.67473.
\]
\[
t = \frac{0.67473}{1.55136 \times 10^{-10}\ \mathrm{yr^{-1}}} = \mathbf{4.349 \times 10^9\ \mathrm{years}} = \mathbf{4.35\ \mathrm{Ga}}.
\]

\textbf{3. Geological Context:}
The Jack Hills zircon crystallized $\mathbf{4.35\ \mathrm{billion\ years}}$ ago, merely $190\ \mathrm{million\ years}$ after the formation of the Earth ($4.54\ \mathrm{Ga}$). This confirms that solid continental crust and liquid water oceans existed astonishingly early in Earth's history!
\end{solution}

% ===========================================================================
% EXERCISE 3.18
% ===========================================================================
\begin{exercise}[$\bigstar\bigstar\bigstar$ Advanced]{Branching Decay Kinetics and Potassium-Argon Dating}
Potassium-40 ($\ce{^{40}_{19}K}$) is an unstable primordial isotope that decays via two competing parallel pathways:
\begin{align*}
\ce{^{40}_{19}K + ^{0}_{-1}e &-> ^{40}_{18}Ar + \nu_e} \quad (\text{Electron Capture, Branching Ratio } B_{\text{EC}} = 0.1072) \\
\ce{^{40}_{19}K &-> ^{40}_{20}Ca + ^{0}_{-1}e + \bar{\nu}_e} \quad (\beta^-\text{ Decay, Branching Ratio } B_{\beta} = 0.8928)
\end{align*}
The total half-life is $t_{1/2} = 1.248 \times 10^9\ \mathrm{years}$ ($\lambda = \lambda_{\text{EC}} + \lambda_\beta$).
\begin{enumerate}[leftmargin=1.5em,itemsep=2pt]
  \item Derive the equation relating accumulated radiogenic Argon-40 to remaining Potassium-40:
  \[
  N_{\ce{Ar}}(t) = \frac{\lambda_{\text{EC}}}{\lambda} N_{\ce{K}}(t) \left( e^{\lambda t} - 1 \right).
  \]
  \item A volcanic basalt sample erupted during an ancient rift event contains $m_{\ce{K}} = 2.500\ \mu\mathrm{g}$ of $\ce{^{40}K}$ and $V_{\ce{Ar}} = 1.250 \times 10^{-5}\ \mathrm{cm^3}$ of radiogenic $\ce{^{40}Ar}$ gas at STP ($V_0 = 22\,414\ \mathrm{cm^3\cdot mol^{-1}}$).
  Calculate the crystallization age of the volcanic rock.
\end{enumerate}
\end{exercise}

\begin{tabletworkspace}[5.0cm]
\end{tabletworkspace}

\begin{solution}
\textbf{1. Derivation for Branching Decay:}
The total rate of loss of $\ce{^{40}K}$ is $\frac{\dd N_{\ce{K}}}{\dd t} = -\lambda N_{\ce{K}}$, where $\lambda = \lambda_{\text{EC}} + \lambda_\beta$.
The rate of generation of Argon-40 is exclusively through electron capture:
\[
\frac{\dd N_{\ce{Ar}}}{\dd t} = \lambda_{\text{EC}} N_{\ce{K}}(t) = \lambda_{\text{EC}} N_0 e^{-\lambda t}.
\]
Integrate from $0$ to $t$ (assuming initial molten rock outgassed all pre-existing argon, $N_{\ce{Ar}}(0) = 0$):
\[
N_{\ce{Ar}}(t) = \lambda_{\text{EC}} N_0 \int_0^t e^{-\lambda t'} \dd t' = \frac{\lambda_{\text{EC}}}{\lambda} N_0 \left( 1 - e^{-\lambda t} \right).
\]
Substitute $N_0 = N_{\ce{K}}(t) e^{\lambda t}$:
\[
\mathbf{N_{\ce{Ar}}(t) = \frac{\lambda_{\text{EC}}}{\lambda} N_{\ce{K}}(t) \left( e^{\lambda t} - 1 \right)}. \quad \text{Q.E.D.}
\]

\textbf{2. Numerical Calculation:}
Total decay constant:
\[
\lambda = \frac{\ln 2}{1.248 \times 10^9\ \mathrm{yr}} = 5.554 \times 10^{-10}\ \mathrm{yr^{-1}}.
\]
Branching fraction: $\frac{\lambda_{\text{EC}}}{\lambda} = 0.1072$.
Number of $\ce{^{40}K}$ nuclei:
\[
N_{\ce{K}} = \frac{2.500 \times 10^{-6}\ \mathrm{g}}{40.0\ \mathrm{g\cdot mol^{-1}}} \times 6.022 \times 10^{23} = 3.7638 \times 10^{16}\ \text{nuclei}.
\]
Number of $\ce{^{40}Ar}$ atoms:
\[
N_{\ce{Ar}} = \frac{1.250 \times 10^{-5}\ \mathrm{cm^3}}{22\,414\ \mathrm{cm^3\cdot mol^{-1}}} \times 6.022 \times 10^{23} = 3.3584 \times 10^{14}\ \text{atoms}.
\]
The ratio is:
\[
\frac{N_{\ce{Ar}}}{N_{\ce{K}}} = \frac{3.3584 \times 10^{14}}{3.7638 \times 10^{16}} = 8.9229 \times 10^{-3}.
\]
From the derived equation:
\[
e^{\lambda t} - 1 = \frac{1}{0.1072} \times \left( 8.9229 \times 10^{-3} \right) = 0.083236 \implies e^{\lambda t} = 1.083236.
\]
\[
\lambda t = \ln(1.083236) = 0.07995.
\]
\[
t = \frac{0.07995}{5.554 \times 10^{-10}\ \mathrm{yr^{-1}}} = \mathbf{1.440 \times 10^8\ \mathrm{years}} = \mathbf{144\ \mathrm{Ma}} \quad \text{(Late Jurassic period)}.
\]
\end{solution}

% ===========================================================================
% EXERCISE 3.19
% ===========================================================================
\begin{exercise}[$\bigstar\bigstar\bigstar$ Advanced]{Cumulative Radiation Dose and Biological Half-Life}
In targeted thyroid radiotherapy, a patient receives an oral capsule containing $A_0 = 3.700\ \mathrm{GBq}$ ($100.0\ \mathrm{mCi}$) of Iodine-131. The isotope is eliminated from the body simultaneously by radioactive decay ($\lambda_p$, physical half-life $t_{p} = 8.020\ \mathrm{days}$) and by biological excretion in urine and sweat ($\lambda_b$, biological clearance half-life $t_{b} = 21.00\ \mathrm{days}$).
\begin{enumerate}[leftmargin=1.5em,itemsep=2pt]
  \item Derive the \textbf{effective half-life} formula:
  \[
  \frac{1}{t_{\text{eff}}} = \frac{1}{t_p} + \frac{1}{t_b} \implies t_{\text{eff}} = \frac{t_p \cdot t_b}{t_p + t_b}.
  \]
  \item Calculate the effective half-life $t_{\text{eff}}$ of Iodine-131 in this patient in days.
  \item The total number of disintegrations occurring inside the patient's body over all time is $N_{\text{cum}} = \int_0^\infty A_{\text{in vivo}}(t) \dd t$. Prove that $N_{\text{cum}} = \dfrac{A_0}{\lambda_{\text{eff}}}$, and compute its numerical value.
\end{enumerate}
\end{exercise}

\begin{tabletworkspace}[5.0cm]
\end{tabletworkspace}

\begin{solution}
\textbf{1. Effective Half-Life Derivation:}
The in vivo disappearance of radioiodine follows two independent parallel first-order rate processes:
\[
\frac{\dd N}{\dd t} = -(\lambda_p + \lambda_b) N(t) \implies \lambda_{\text{eff}} = \lambda_p + \lambda_b.
\]
Expressing rate constants in terms of half-lives ($\lambda = \ln 2 / t_{1/2}$):
\[
\frac{\ln 2}{t_{\text{eff}}} = \frac{\ln 2}{t_p} + \frac{\ln 2}{t_b} \implies \mathbf{\frac{1}{t_{\text{eff}}} = \frac{1}{t_p} + \frac{1}{t_b}} \implies \mathbf{t_{\text{eff}} = \frac{t_p \cdot t_b}{t_p + t_b}}. \quad \text{Q.E.D.}
\]

\textbf{2. Numerical Calculation of $t_{\text{eff}}$:}
Given $t_p = 8.020\ \mathrm{days}$ and $t_b = 21.00\ \mathrm{days}$:
\[
t_{\text{eff}} = \frac{8.020 \times 21.00}{8.020 + 21.00} = \frac{168.42}{29.02} = \mathbf{5.804\ \mathrm{days}}.
\]

\textbf{3. Cumulative Total In Vivo Disintegrations:}
The in vivo activity is $A_{\text{in vivo}}(t) = A_0 e^{-\lambda_{\text{eff}} t}$.
Integrating over $t \in [0, \infty)$:
\[
N_{\text{cum}} = \int_0^\infty A_0 e^{-\lambda_{\text{eff}} t} \dd t = A_0 \left[ -\frac{e^{-\lambda_{\text{eff}} t}}{\lambda_{\text{eff}}} \right]_0^\infty = A_0 \left( 0 - \left( -\frac{1}{\lambda_{\text{eff}}} \right) \right) = \mathbf{\frac{A_0}{\lambda_{\text{eff}}}}. \quad \text{Q.E.D.}
\]
Convert $t_{\text{eff}}$ to seconds:
\[
t_{\text{eff}} = 5.804 \times 86\,400\ \mathrm{s} = 5.0147 \times 10^5\ \mathrm{s}.
\]
\[
\lambda_{\text{eff}} = \frac{\ln 2}{5.0147 \times 10^5\ \mathrm{s}} = 1.3823 \times 10^{-6}\ \mathrm{s^{-1}}.
\]
Given $A_0 = 3.700 \times 10^9\ \mathrm{Bq}$:
\[
N_{\text{cum}} = \frac{3.700 \times 10^9\ \mathrm{s^{-1}}}{1.3823 \times 10^{-6}\ \mathrm{s^{-1}}} = \mathbf{2.677 \times 10^{15}\ \text{nuclear disintegrations}}.
\]
\end{solution}

% ===========================================================================
% EXERCISE 3.20
% ===========================================================================
\begin{exercise}[$\bigstar\bigstar\bigstar$ Advanced]{The Bateman Differential Equations for Decay Cascades}
Consider a sequential radioactive decay chain with parent nucleus $1$ decaying into daughter $2$, which in turn decays into stable granddaughter $3$:
\[
\ce{1} \xrightarrow{\lambda_1} \ce{2} \xrightarrow{\lambda_2} \ce{3 (stable)}.
\]
At time $t = 0$, only pure parent is present: $N_1(0) = N_0$ and $N_2(0) = 0$.
\begin{enumerate}[leftmargin=1.5em,itemsep=2pt]
  \item Write the coupled system of differential rate equations for $\dd N_1 / \dd t$ and $\dd N_2 / \dd t$.
  \item Solve the system analytically using an integrating factor to prove the Bateman formula:
  \[
  N_2(t) = \frac{\lambda_1}{\lambda_2 - \lambda_1} N_0 \left( e^{-\lambda_1 t} - e^{-\lambda_2 t} \right).
  \]
  \item Find the exact time $t_{\max}$ at which the daughter activity $A_2(t) = \lambda_2 N_2(t)$ reaches its maximum.
  \item Apply this to the medical isotope generator system Molybdenum-99 ($\lambda_1 = 0.01050\ \mathrm{h^{-1}}$, $t_{1/2,1} = 66.0\ \mathrm{h}$) producing Technetium-99m ($\lambda_2 = 0.11533\ \mathrm{h^{-1}}$, $t_{1/2,2} = 6.01\ \mathrm{h}$). Calculate $t_{\max}$ in hours.
\end{enumerate}
\end{exercise}

\begin{tabletworkspace}[5.0cm]
\end{tabletworkspace}

\begin{solution}
\textbf{1. Coupled Differential Equations:}
\begin{align}
\frac{\dd N_1}{\dd t} &= -\lambda_1 N_1 \label{eq:b1} \\
\frac{\dd N_2}{\dd t} &= \lambda_1 N_1 - \lambda_2 N_2 \label{eq:b2}
\end{align}

\textbf{2. Analytical Solution by Integrating Factor:}
From \eqref{eq:b1}, $N_1(t) = N_0 e^{-\lambda_1 t}$. Substitute into \eqref{eq:b2}:
\[
\frac{\dd N_2}{\dd t} + \lambda_2 N_2 = \lambda_1 N_0 e^{-\lambda_1 t}.
\]
Multiply by the integrating factor $e^{\lambda_2 t}$:
\[
\frac{\dd}{\dd t} \left( N_2 e^{\lambda_2 t} \right) = \lambda_1 N_0 e^{(\lambda_2 - \lambda_1) t}.
\]
Integrate both sides from $0$ to $t$:
\[
N_2(t) e^{\lambda_2 t} - 0 = \lambda_1 N_0 \left[ \frac{e^{(\lambda_2 - \lambda_1) t} - 1}{\lambda_2 - \lambda_1} \right].
\]
Multiplying across by $e^{-\lambda_2 t}$:
\[
\mathbf{N_2(t) = \frac{\lambda_1}{\lambda_2 - \lambda_1} N_0 \left( e^{-\lambda_1 t} - e^{-\lambda_2 t} \right)}. \quad \text{Q.E.D.}
\]

\textbf{3. Determination of Maximum Time $t_{\max}$:}
Set the derivative $\frac{\dd N_2}{\dd t} = 0$:
\[
-\lambda_1 e^{-\lambda_1 t_{\max}} + \lambda_2 e^{-\lambda_2 t_{\max}} = 0 \implies \lambda_1 e^{-\lambda_1 t_{\max}} = \lambda_2 e^{-\lambda_2 t_{\max}}.
\]
\[
\frac{e^{-\lambda_1 t_{\max}}}{e^{-\lambda_2 t_{\max}}} = \frac{\lambda_2}{\lambda_1} \implies e^{(\lambda_2 - \lambda_1) t_{\max}} = \frac{\lambda_2}{\lambda_1}.
\]
Taking natural logarithms:
\[
(\lambda_2 - \lambda_1) t_{\max} = \ln \left( \frac{\lambda_2}{\lambda_1} \right) \implies \mathbf{t_{\max} = \frac{\ln(\lambda_2 / \lambda_1)}{\lambda_2 - \lambda_1}}.
\]

\textbf{4. Application to $\ce{^{99}Mo} \to \ce{^{99m}Tc}$ Generator:}
Given $\lambda_1 = 0.01050\ \mathrm{h^{-1}}$ and $\lambda_2 = 0.11533\ \mathrm{h^{-1}}$:
\[
\lambda_2 - \lambda_1 = 0.11533 - 0.01050 = 0.10483\ \mathrm{h^{-1}}.
\]
\[
\frac{\lambda_2}{\lambda_1} = \frac{0.11533}{0.01050} = 10.9838 \implies \ln(10.9838) = 2.3964.
\]
\[
t_{\max} = \frac{2.3964}{0.10483\ \mathrm{h^{-1}}} = \mathbf{22.86\ \mathrm{hours}} \approx \mathbf{23\ \mathrm{hours}}.
\]
In hospital radiopharmacy practice, "milking" the molybdenum "cow" (technetium generator) is performed every 24 hours precisely because daughter technetium activity reaches its peak approximately 23 hours after each elution!
\end{solution}
